Extending all-weather road access to dispersed rural hamlets is a major source of new linear infrastructure and of land disturbance, yet access routes are usually planned house by house from manually digitised building inventories. We present a framework that links satellite building detection to a cost-surface-driven shared road network and apply it to Damarakunta, an agricultural village of 18.5 square kilometres in Siddipet district, Telangana, India. Buildings were extracted from 0.6 metre imagery with a fully fine-tuned self-supervised vision-transformer segmentation network, which reached 53.8 per cent mean intersection over union (68.9 per cent mean Dice, 71.3 per cent pixel accuracy) on a public land-cover benchmark and agreed with 71.5 per cent of village candidates in an independent building layer. A ten-metre cost surface built from land cover, slope, hard terrain constraints and a penalty on building footprints drives a two-stage planner in which routes may end on existing roads or on other planned routes. For 163 residential targets and five facilities, all 168 targets were connected with no constraint violation. The shared network needs 5,069 metres of new corridor, 24.9 per cent (1,684 metres) less than independent routing, about 0.84 hectares less land under an assumed five-metre formation. The footprint penalty halves the routes crossing other dwellings, from 56 to 28, for a 7.7 per cent longer network. Across 23 simulated omission and addition runs the saving stayed between 24.9 and 26.6 per cent, but omitted buildings shortened the plan by up to 29.4 per cent and silently removed up to 38.0 per cent of households from it. Shared planning can therefore lower the land footprint of rural roads, whereas incomplete building inventories bias land-take estimates downward and should be reported as ranges. Results are model outputs for one village without field validation.
